\chapter{Conclusions}

When this project started, the problem seemed straightforward: build a tool that lets a small group of students share files on a local network without setting up servers or learning Git. The research and design process revealed that the problem is more nuanced than it appears. Distributed synchronization, even for a handful of peers, touches on event ordering, consistency models, and conflict resolution, all topics with decades of research behind them.

The literature review (Chapter~2) showed that existing tools each solve part of the problem. Cloud services handle synchronization well but depend on external infrastructure. Git handles versioning well but requires technical expertise. Syncthing handles peer-to-peer transfer well but treats versioning as an afterthought. No existing tool scores well across all five evaluation dimensions (membership, synchronization, history, file support, and transparency) simultaneously.

The system described in this thesis fills that gap. The Go networking daemon discovers peers through mDNS/DNS-SD and coordinates synchronization using per-file Lamport clocks with Last-Write-Wins conflict resolution. The C++ storage daemon watches the filesystem for changes, computes content hashes, and creates backups before every overwrite. The Java frontend shows users what is happening without requiring them to understand the underlying protocols.

The design makes three deliberate trade-offs:

\begin{itemize}
    \item \textbf{Lamport clocks instead of vector clocks.} Simpler to implement and manage when peers join and leave, but unable to detect truly concurrent edits.
    \item \textbf{LWW instead of CRDTs.} One version is chosen automatically; the other is preserved in a backup. Simple and predictable, but not conflict-free in the CRDT sense.
    \item \textbf{Whole-file transfers instead of delta sync.} Easy to implement and debug, but wasteful for large files with small changes.
\end{itemize}

Each of these trade-offs is appropriate for 3 to 5 peers on a local network with a project timeline of a few months. A production system serving hundreds of users would need to revisit all three.

Several directions for future work remain:

\begin{itemize}
    \item \textbf{Block-level delta synchronization.} Implementing the rsync rolling-checksum algorithm in the C++ daemon to transfer only changed blocks, reducing bandwidth for large files.
    \item \textbf{CRDT-based metadata.} Extending the SQLite schema to use state-based CRDTs for file metadata, enabling conflict-free convergence without LWW.
    \item \textbf{Cross-subnet discovery.} Supporting peers on different subnets, possibly through a lightweight relay, while preserving zero-configuration on local networks.
    \item \textbf{Mobile support.} Adapting the architecture for Android and iOS, where background execution and filesystem access are more restricted than on desktop operating systems.
\end{itemize}

What this project demonstrates is that the individual building blocks for lightweight peer-to-peer synchronization, mDNS, Lamport clocks, LWW, TCP, SHA-256, are well-understood and well-documented. The challenge is not inventing new algorithms. It is combining existing ones into a package that is simple enough for non-technical users, safe enough that no data is lost, and light enough to run unnoticed on a student laptop. That is what this system sets out to do.
